\chapter{Méthodologie de constitution des données}

\section{Sources de données}

Contrairement au périmètre initialement envisagé (données simulées calibrées sur documents
de référence), ce travail s'appuie in fine sur \textbf{neuf dossiers de crédit réels}, extraits
manuellement des captures d'écran de l'outil interne de notation MNS2 du Centre d'Affaires
Al Istiqlal. Ces dossiers couvrent trois segments de crédit distincts :

\begin{itemize}
    \item \textbf{Crédit de Fonctionnement / Court Terme (CT)} : 4 dossiers (Feuil2, Feuil3, Feuil8, Feuil5) ;
    \item \textbf{Promotion Immobilière (PI)} : 3 dossiers (Feuil1, Feuil13, Feuil4) ;
    \item \textbf{Moyen/Long Terme (MLT)} : 2 dossiers dont le segment reste à confirmer (Feuil11, Feuil12).
\end{itemize}

\begin{limite}[Taille de l'échantillon]
Neuf dossiers constituent un échantillon extrêmement restreint au regard des standards de la
modélisation statistique. Ce choix n'est pas une préférence méthodologique mais une contrainte
opérationnelle : ces dossiers représentent l'ensemble des cas traités et documentés durant la
période de stage. Toutes les analyses qui suivent doivent être lues à la lumière de cette
contrainte, explicitée dès le Chapitre 1.
\end{limite}

\section{Fusion et harmonisation des variables}

Chaque dossier ayant été saisi indépendamment à partir de captures d'écran, la première étape
a consisté à fusionner les neuf tableurs sources en une base unique. Cette fusion brute a fait
apparaître \textbf{638 colonnes de variables}, un nombre très supérieur au nombre de variables
réellement distinctes, en raison de :

\begin{enumerate}
    \item \textbf{Variantes orthographiques} du même intitulé de variable (ex. fautes de frappe
    lors de la saisie manuelle) ;
    \item \textbf{Doublons structurels} : certaines variables apparaissent plusieurs fois sous
    des suffixes \texttt{\_\_dup2}, \texttt{\_\_dup3}, vraisemblablement liés à des exercices
    comptables différents (N-2, N-1, N) mal étiquetés à la source ;
    \item \textbf{Confusions actif/passif} : trois groupes de variables mêlaient des postes
    d'actif et de passif sous un intitulé commun (ex. \texttt{tresorerie\_actif} et
    \texttt{tresorerie\_passif}), nécessitant une séparation manuelle stricte pour ne pas
    fusionner des grandeurs comptables non comparables.
\end{enumerate}

Un dictionnaire de correspondance (\texttt{variable brute} $\rightarrow$ \texttt{variable
canonique}) a été construit et appliqué par coalescence : pour chaque dossier et chaque groupe
de variables équivalentes, la première valeur non nulle rencontrée est retenue. Cette opération
réduit la base de \textbf{638 à 339 variables canoniques}, tout en conservant une traçabilité
complète du mapping appliqué (table \texttt{mapping\_nettoyage} de la base de données, cf.
Section~3.4).

\begin{limite}
La coalescence par « première valeur non nulle » constitue une hypothèse simplificatrice.
Si les colonnes \texttt{\_\_dup2}/\texttt{\_\_dup3} correspondent effectivement à des exercices
comptables différents plutôt qu'à des doublons de saisie, cette opération entraîne une perte
d'information (seul un exercice est conservé par variable). Ce point est identifié comme
prioritaire pour validation avec le maître de stage.
\end{limite}

Une correction d'échelle a par ailleurs été appliquée sur la variable \texttt{cash\_flow} : un
dossier présentait une valeur exprimée en dirhams (DH) alors que les autres étaient en milliers
de dirhams (KDH), introduisant un facteur d'échelle de 1000 corrigé après détection automatique
des incohérences d'unité.

Après nettoyage, la couverture des variables reste très hétérogène : sur les 339 variables
canoniques, seules 55 sont renseignées pour au moins 6 des 9 dossiers. Cette forte sparsité a
directement orienté le choix du schéma de base de données (Section~3.4) et la sélection des
variables retenues pour la modélisation (Chapitre~5).

\section{Anonymisation et sécurisation des données}

Une exigence de sécurité et de confidentialité a été posée dès la conception du système :
\textbf{aucune donnée identifiante (raison sociale) ne doit transiter par le moteur de scoring,
le moteur de décision, ni les modules d'explicabilité.} Concrètement :

\begin{itemize}
    \item les raisons sociales sont isolées dans une table dédiée (\texttt{identites}), distincte
    de la table des dossiers (\texttt{dossiers}), reliées uniquement par un identifiant technique
    anonyme (\texttt{dossier\_id}) ;
    \item les modules de scoring, de machine learning et de feature importance
    (Chapitres~4 et~5) ne référencent \emph{jamais} la table \texttt{identites} ;
    \item dans l'application (Chapitre~6), seule la vue « Direction » — habilitée à statuer sur
    le dossier — effectue la jointure permettant de lever l'anonymat.
\end{itemize}

Cette séparation répond à un principe de \emph{minimisation des données} : chaque composant du
système n'accède qu'aux informations strictement nécessaires à sa fonction.

\section{Construction de la base de données}

La base de données finale (SQLite, embarquée, sans dépendance réseau) est structurée autour de
sept tables principales, résumées au Tableau~\ref{tab:schema-bdd}.

\begin{table}[h]
\centering
\small
\begin{tabular}{@{}p{3.2cm}p{9.5cm}@{}}
\toprule
\textbf{Table} & \textbf{Rôle} \\
\midrule
\texttt{dossiers} & Une ligne par dossier (identifiant, segment) — anonyme \\
\texttt{identites} & Correspondance dossier\_id $\rightarrow$ raison sociale — isolée \\
\texttt{variables\_dictionnaire} & Catalogue des 339 variables canoniques (couverture, type) \\
\texttt{variables\_valeurs} & Valeurs au format \emph{entité-attribut-valeur} (EAV) \\
\texttt{mapping\_nettoyage} & Traçabilité variable brute $\rightarrow$ variable canonique \\
\texttt{notations} & Notation MNS2 réelle reconstituée par dossier \\
\texttt{scores\_calcules} / \texttt{decision\_ml} & Sorties des systèmes de scoring et de décision (Ch.~4-6) \\
\bottomrule
\end{tabular}
\caption{Schéma simplifié de la base de données}
\label{tab:schema-bdd}
\end{table}

Le choix d'un modèle \textbf{entité-attribut-valeur (EAV)} pour la table
\texttt{variables\_valeurs} — plutôt qu'une table large de 339 colonnes — est directement motivé
par la sparsité constatée en Section~3.2 : une table large contiendrait, pour la majorité des
variables, plus de 30\,\% de cellules vides, ce qui est inefficace en stockage et rend les
requêtes d'agrégation plus complexes. Le format EAV ne stocke que les couples
(dossier, variable, valeur) effectivement renseignés.

\section{Synthèse}

Cette phase de constitution des données, bien qu'en amont du cœur méthodologique du projet
(scoring et décision), a représenté une part substantielle du travail réalisé : elle conditionne
directement la fiabilité de toutes les analyses ultérieures. Les Chapitres~4 et~5 s'appuient sur
cette base nettoyée et anonymisée pour construire, respectivement, le score expert et le score
par apprentissage automatique.
